\chapter{État de l'art et méthodologie}
\label{chap:etatart}

Ce chapitre situe ce travail par rapport aux approches existantes de
détection d'anomalies dans les séries temporelles de trafic réseau,
puis détaille la méthodologie retenue, du prétraitement des données
brutes jusqu'au calcul du score d'anomalie.

\clearpage
\section{État de l'art}

\subsection{Approches statistiques et par seuils}

Les premières approches de détection d'anomalies dans des séries
temporelles reposent sur des méthodes statistiques classiques :
cartes de contrôle, moyennes mobiles, ou modèles de prévision (par
exemple ARIMA) dont on surveille le résidu -- l'écart entre la valeur
observée et la valeur prédite. Une anomalie est alors signalée
lorsque ce résidu, ou plus simplement la valeur elle-même, dépasse un
seuil fixé a priori ou calibré statistiquement (par exemple à partir
de l'écart type de la série). Ces méthodes sont simples à mettre en
œuvre et à interpréter, mais elles supposent une certaine stationnarité
du signal et se révèlent sensibles au choix du seuil, ainsi qu'aux
variations normales (saisonnalité, effets de calendrier) qui peuvent
être confondues avec de véritables anomalies \parencite{chandola2009anomaly}.

\subsection{Approches par apprentissage automatique}

Les approches par apprentissage automatique cherchent à apprendre
directement, à partir des données, ce qu'est un comportement
« normal », sans reposer sur un seuil fixé manuellement. Les méthodes
non supervisées -- clustering, forêts d'isolement, auto-encodeurs --
identifient les observations qui s'écartent le plus de la structure
apprise sur l'ensemble des données, tandis que les approches par
apprentissage profond (réseaux récurrents, auto-encodeurs séquentiels)
permettent de capturer des dépendances temporelles plus complexes
\parencite{chalapathy2019deep}. Ces méthodes sont en général plus
coûteuses à mettre en œuvre et moins directement interprétables, et
nécessitent souvent un volume de données d'entraînement important
pour être fiables.

\subsection{Approches par encodage de signatures et distance d'édition}

Une troisième famille d'approches encode le comportement d'une entité
(ici, le trafic horaire d'une zone géographique) sous la forme d'une
séquence symbolique -- une « signature » -- puis mesure l'écart entre
deux séquences à l'aide d'une distance d'édition, empruntée à
l'algorithmique des chaînes de caractères. C'est l'approche retenue
dans le framework DiNATrAX~\parencite{maudoux2024dinatrax}, qui encode
l'activité d'une entité réseau en une séquence de lettres et détecte
les anomalies par distance de Damerau-Levenshtein à un profil de
référence. Cette approche présente trois avantages pour le cas
étudié~: elle ne nécessite pas de données étiquetées, elle reste
interprétable (chaque lettre correspond directement à un niveau de
trafic observable), et elle est sensible aux variations d'amplitude du
signal -- contrairement à un découpage en quantiles, qui absorbe une
hausse uniforme du trafic sans faire apparaître de changement de
signature. C'est cette approche qui a été adaptée et retenue pour ce
travail.

\clearpage
\section{Méthodologie}

\subsection{Vue d'ensemble du pipeline}

Le traitement se décompose en six étapes principales, illustrées par
la figure~\ref{fig:pipeline} : agrégation du trafic brut, encodage en
signature journalière, construction d'un profil de référence,
calcul de la distance d'édition, conversion en score z, puis
recoupement calendaire des anomalies détectées. Ce pipeline est
exécuté indépendamment sur chacun des trois canaux disponibles
(Internet, SMS, appels).

\begin{figure}[htbp]
  \centering
  \begin{tikzpicture}[
    node distance=6mm,
    box/.style={rectangle, draw, rounded corners, fill=parisneutral!12,
                text width=8.4cm, align=center, minimum height=9mm, font=\small},
    arr/.style={-{Stealth[length=2mm]}, thick}
  ]
    \node[box] (a) {Trafic brut par (grille, jour, heure) -- 62 fichiers, \textasciitilde{}4M lignes chacun};
    \node[box, below=of a] (b) {Encodage en signature de 24 lettres (A-F, multiples de la médiane horaire de la grille)};
    \node[box, below=of b] (c) {Profil de référence par (grille, jour de semaine) -- consensus lettre à lettre};
    \node[box, below=of c] (d) {Distance de Damerau-Levenshtein (signature du jour $\leftrightarrow$ profil de référence)};
    \node[box, below=of d] (e) {Score z par grille -- jour anormal si $z > 3$};
    \node[box, below=of e] (f) {Recoupement avec le calendrier d'événements (jours fériés, matchs San Siro)};
    \draw[arr] (a) -- (b);
    \draw[arr] (b) -- (c);
    \draw[arr] (c) -- (d);
    \draw[arr] (d) -- (e);
    \draw[arr] (e) -- (f);
  \end{tikzpicture}
  \caption{Vue d'ensemble du pipeline de détection}
  \label{fig:pipeline}
\end{figure}

\subsection{Agrégation du trafic}

Les fichiers sources contiennent, pour chaque interaction (connexion
Internet, SMS, appel), un horodatage et l'identifiant de la grille
géographique parmi les 10 000 zones carrées composant le maillage de
Milan. Ces événements sont agrégés par \emph{(grille, jour, heure)},
ce qui produit, pour chaque grille et chaque canal, une série
temporelle horaire de trafic sur toute la période couverte.

\subsection{Encodage en signature journalière}

Pour chaque grille, le trafic horaire est encodé en un niveau discret
parmi six (tableau~\ref{tab:levels}), exprimé en multiple de la
médiane horaire \emph{propre à cette grille}. Ce choix rend l'échelle
relative -- une grille calme et une grille du centre-ville restent
comparables -- tout en la rendant sensible à l'amplitude du trafic~:
un doublement uniforme du trafic d'une grille fait basculer ses
lettres vers des niveaux supérieurs, ce qu'un découpage en quantiles
ne peut pas détecter. La signature d'une journée est la concaténation
des 24 lettres, une par heure.

\begin{table}[htbp]
  \centering
  \caption{Niveaux de la signature, en multiple de la médiane horaire de la grille}
  \label{tab:levels}
  \begin{tabular}{ccl}
    \toprule
    \textbf{Lettre} & \textbf{Multiple de la médiane} & \textbf{Signification} \\
    \midrule
    A & $< 0{,}43\times$          & Activité très faible -- creux profond \\
    B & $0{,}43\times$ -- $0{,}63\times$ & Activité faible \\
    C & $0{,}63\times$ -- $1{,}0\times$  & En dessous du niveau habituel \\
    D & $1{,}0\times$ -- $1{,}35\times$  & Niveau habituel \\
    E & $1{,}35\times$ -- $2{,}1\times$  & Activité soutenue \\
    F & $> 2{,}1\times$           & Pointe d'activité \\
    \bottomrule
  \end{tabular}
\end{table}

Ces seuils ont été calibrés sur les percentiles réels du canal
Internet (1/5/20/50/80/95/99) plutôt que fixés à des valeurs rondes~:
cette calibration réduit le nombre de journées détectées comme
anormales de moitié par rapport à des seuils ronds (0,48~\% contre
0,98~\% des journées), pour un taux de correspondance calendaire plus
élevé. Les canaux SMS et appels, dont les dynamiques diffèrent
sensiblement du canal Internet, sont calibrés séparément selon le même
principe.

\subsection{Construction du profil de référence}

Pour chaque grille, un profil de référence est construit \emph{par
jour de semaine} (tous les mardis, tous les mercredis, etc.)~: à
chaque heure, on retient la lettre majoritaire observée parmi les
journées de ce jour de semaine. Ce profil consensus constitue
l'équivalent d'un comportement « normal » propre à chaque grille et à
chaque jour de semaine.

\subsection{Distance de Damerau-Levenshtein}

L'écart entre la signature d'une journée réelle et son profil de
référence est mesuré par la distance de Damerau-Levenshtein, qui
compte le nombre minimal d'opérations élémentaires -- insertion,
suppression, substitution, ou transposition de deux lettres adjacentes
-- nécessaires pour transformer une séquence en l'autre. L'ajout de la
transposition par rapport à la distance de Levenshtein classique est
important ici~: un décalage d'une heure dans un pic de trafic se
traduit par la permutation de deux lettres adjacentes, que la
distance de Levenshtein pénaliserait à tort comme deux erreurs
distinctes.

Le tableau~\ref{tab:dl-example} illustre ce calcul sur un exemple
réduit à quatre lettres, en comparant la séquence \texttt{BCDE} à une
référence \texttt{BDCE} -- les lettres \texttt{C} et \texttt{D} y sont
simplement permutées. La distance de Damerau-Levenshtein vaut alors 1
(une transposition), alors qu'une distance de Levenshtein classique
compterait 2 (deux substitutions).

\begin{table}[htbp]
  \centering
  \caption{Exemple de calcul de la distance de Damerau-Levenshtein entre \texttt{BCDE} et \texttt{BDCE}}
  \label{tab:dl-example}
  \begin{tabular}{c|ccccc}
    \toprule
     & \textbf{\o} & \textbf{B} & \textbf{D} & \textbf{C} & \textbf{E} \\
    \midrule
    \textbf{\o} & 0 & 1 & 2 & 3 & 4 \\
    \textbf{B}  & 1 & 0 & 1 & 2 & 3 \\
    \textbf{C}  & 2 & 1 & 1 & 1 & 2 \\
    \textbf{D}  & 3 & 2 & 1 & \textbf{1} & 2 \\
    \textbf{E}  & 4 & 3 & 2 & 2 & \textbf{1} \\
    \bottomrule
  \end{tabular}
\end{table}

\subsection{Score z et seuil de détection}

Une distance élevée en valeur absolue n'est pas directement
interprétable~: certaines grilles ont naturellement des journées plus
variables que d'autres. Les distances sont donc normalisées en score
z \emph{par grille} -- c'est-à-dire centrées et réduites par la
moyenne et l'écart type des distances observées sur cette grille --
puis comparées à un seuil fixe. Une journée est marquée comme
anormale lorsque son score z dépasse 3 (figure~\ref{fig:zscore}),
c'est-à-dire lorsqu'elle s'écarte de plus de trois écarts types du
comportement habituel de sa grille.

\begin{figure}[htbp]
  \centering
  \begin{tikzpicture}
    \begin{axis}[
      width=0.68\textwidth, height=4.6cm,
      domain=-4:4, samples=200,
      axis lines=middle,
      xlabel={score $z$}, ylabel={densité},
      xtick={-3,-2,-1,0,1,2,3}, ytick=\empty,
      every axis x label/.style={at={(current axis.right of origin)}, anchor=west},
      every axis y label/.style={at={(current axis.above origin)}, anchor=south},
      clip=false,
    ]
      \addplot[thick, black] {1/sqrt(2*pi)*exp(-x^2/2)};
      \addplot[draw=none, fill=parisneutral, fill opacity=0.45, domain=3:4]
        {1/sqrt(2*pi)*exp(-x^2/2)} \closedcycle;
      \node[font=\footnotesize] at (axis cs:3.55,0.09) {jour anormal};
    \end{axis}
  \end{tikzpicture}
  \caption{Zone d'anomalie ($z > 3$) sur la distribution des scores z}
  \label{fig:zscore}
\end{figure}

\subsection{Détection multi-canal et recoupement calendaire}

Ce pipeline est exécuté indépendamment sur les trois canaux
disponibles (Internet, SMS, appels), chacun avec ses propres seuils
de nomenclature. Une anomalie confirmée simultanément par au moins
deux canaux est considérée comme plus fiable qu'une anomalie détectée
sur un seul canal. Enfin, les anomalies détectées sont recoupées avec
un calendrier construit manuellement (jours fériés italiens, matchs à
domicile de l'AC Milan et de l'Inter à San Siro), afin de distinguer
une anomalie expliquée par un événement connu d'un signal réellement
inattendu. Les résultats de cette validation sont présentés au
chapitre~\ref{chap:conception}.
